\subsection{Statement and background}

For a finite or infinite word \(x\), a factor is a contiguous subword.
Write \(\Fac_m(x)\) for its set of length-\(m\) factors. A square is
\(UU\) with \(U\ne\e\), and a word is square-free if it has no square
factor. A factor of a one-sided infinite word is recurrent if its
set of starting positions is infinite. Positions in infinite words
are numbered from one.

\begin{theoremT}\label{thm:T}
Let \(\xi\) be a real algebraic irrational and let \(\mathbf a\) be the
canonical ternary digit word of its fractional part. Write
\[
P=\{\tern{010},\tern{020},\tern{101},\tern{121},\tern{202},\tern{212}\}.
\]
Then at least one of the following holds.
\begin{enumerate}
\item Infinitely many positions of \(\mathbf a\) are starting positions
of a square.
\item At least five distinct members of \(P\) are recurrent in
\(\mathbf a\).
\end{enumerate}
\end{theoremT}

Adamczewski and Rampersad \cite[Theorem 3.1]{AR} proved that an
algebraic ternary expansion has infinitely many square occurrences
or that \(\tern{010}\) or \(\tern{02120}\) is recurrent.
Theorem~T concerns square starting positions and recurrence of
specified words of length three. It gives no unbounded-length conclusion
for square roots or palindromes in Bugeaud's Problem~10.42 \cite{Bugeaud}.

Over a finite alphabet, infinitely many distinct squares are
equivalent to unbounded square-root lengths. Infinite binary words
with finitely many distinct squares exist \cite{BC}. They have
infinitely many square starting positions because every length-four
binary word contains a square.

Equivalently, the base-\(3\) value of a one-sided digit word is
transcendental if the word has a square-free suffix avoiding two
distinct members of \(P\). Indeed, finitely many square starting
positions and finitely many occurrences of the two words can all be
passed by one fixed cut. Such a suffix also excludes eventual
periodicity, so the value is irrational.

We derive Theorem~T as a corollary of established results. The two
one-sided lemmas in Section~\ref{sec:lifting} give the prefix control
needed to use the criterion of Adamczewski, Bugeaud, and Luca
\cite{ABL}.

\subsection{The avoidance inputs}

Put \(\Sigma_3=\{0,1,2\}\) and \(\Sigma_5=\{0,1,2,3,4\}\).
Badkobeh and Ochem \cite[Theorems 1 and 2]{BO} use the morphism
\(f=f_5\) defined by the images
\[
(f(0),f(1),f(2),f(3),f(4))
=(\tern{01},\tern{23},\tern{4},\tern{21},\tern{0}).
\]
Let \(t=f^\omega(0)\) and let
\[
\begin{split}
F_5=\{&
\tern{02},\tern{03},\tern{13},\tern{14},\tern{20},\tern{24},
\tern{31},\tern{32},\tern{40},\tern{41},\tern{43},\\
&\tern{121},\tern{212},\tern{304},\tern{3423},\tern{4234}\}.
\end{split}
\]
Let \(\Y\) consist of all one-sided infinite square-free words on
\(\Sigma_5\) that avoid \(F_5\). Theorem 1 gives \(t\in\Y\). Its proof
supplies the finite input
\begin{equation}\label{eq:window}
\begin{gathered}
pvs\text{ is square-free and avoids }F_5,\qquad
|p|=|s|=32,\quad |v|=8\\
\Longrightarrow\quad v\in\Fac_8(t).
\end{gathered}
\end{equation}
Section~\ref{sec:computation} describes an independent reproduction
of \eqref{eq:window}.

For \(i=1,2\) put
\[
B_1=\{\tern{010},\tern{020}\},\qquad
B_2=\{\tern{121},\tern{212}\}.
\]
Theorem 2 of \cite{BO}, with the definition of characterization
in that paper, asserts that every recurrent factor of any one-sided
infinite square-free word avoiding \(B_i\) occurs in \(s_i=M_i(t)\).
In the order \(0,1,2,3,4\), the images are
\[
M_1=(\tern{012},\tern{1},\tern{02},\tern{12},\e),\qquad
M_2=(\tern{02},\tern{1},\tern{0},\tern{12},\e).
\]
The recurrent languages also appear in Ochem and Rosenfeld
\cite[Theorems 2 and 3]{OR}, where the characterization is expressed
via distances between consecutive occurrences of a letter. Our use
of these classes requires statements about arbitrary one-sided words,
including their initial factors.

\subsection{Two one-sided transfer lemmas}\label{sec:lifting}

Badkobeh and Ochem \cite[\S3.1]{BO} give the four nonempty
images of \(M_i\circ f\), which appear as \(A,B,C,D\) below.
Define \(g_i=M_i\circ f^2\), so that \(s_i=g_i(t)\).
Direct composition gives the five images in the last column.
\[
\begin{array}{c@{\quad}cccc@{\quad}c}
\toprule
i & A & B & C & D & (g_i(0),g_i(1),g_i(2),g_i(3),g_i(4))\\
\midrule
1 & \tern{0121} & \tern{0212} & \tern{012} & \tern{021} & (AB,D,C,B,A)\\
2 & \tern{021} & \tern{012} & \tern{02} & \tern{01} & (AB,D,C,B,A)\\
\bottomrule
\end{array}
\]
Each of \(A,B,C,D\) has its only \(0\) at the beginning.
Every letter image under \(g_i\) has length between \(2\) and \(8\).

\begin{lemma}[Ternary suffix decoding]\label{lem:ternary}
If a one-sided infinite square-free word \(x\) on \(\Sigma_3\)
avoids \(B_i\), then \(x=H g_i(z_0)\) for a finite prefix \(H\)
and \(z_0\in\Y\). More precisely, if every length-\(41\) factor
starting at or after position \(N\) in \(x\) occurs in \(s_i\),
then \(H\) can be chosen with \(|H|\le N+10\).
\end{lemma}

\begin{proof}
Theorem 2 of \cite{BO} ensures that such an \(N\) exists. Each
length-\(41\) word outside \(\Fac_{41}(s_i)\) occurs only finitely
often in \(x\), and there are finitely many such words. Take \(N\)
one greater than the last starting position of any of them, or
take \(N=1\) if there is no such position.

Call the segment from one \(0\) to just before the next \(0\) a return
block. In a square-free ternary word every length-\(4\) factor contains
\(0\). A word of length \(4\) on \(\{1,2\}\) has an adjacent equal
pair or is an alternating square. Thus consecutive zero positions
in \(x\) differ by at most \(4\). If a return block \(R\) starts at
or after \(N\), the length-at-most-\(5\) word \(R0\) occurs in
\(s_i\) by right extension to length \(41\). All zero positions in
\(s_i=g_i(t)\) mark return-block starts by the table. The two zeros
of this occurrence of \(R0\) are consecutive, so \(R\) is one of
\(A,B,C,D\).

In the return-block sequence of \(g_i(t)\), an adjacent \(AB\)
occurrence is exactly the image of \(0\). For such a pair to cross
an image boundary, the image on the left would have to be \(g_i(4)\),
which ends in \(A\), and the one on the right would have to be
\(g_i(3)\), which starts in \(B\). The word \(t\) avoids \(\tern{43}\).
The local rule therefore pairs every adjacent \(AB\) and decodes it
as \(0\). It decodes the unpaired blocks \(D,C,B,A\) as \(1,2,3,4\),
respectively. All pairs are disjoint, and the neighbors of a block
determine whether it belongs to a pair.

Let \(q_0\) be the first zero position in \(x\) at or after \(N\), and
let \(q_1,q_2\) be the next two. Their bounds are \(q_0\le N+3\),
\(q_1\le N+7\), and \(q_2\le N+11\). Apply the pairing rule to the
return-block sequence beginning at \(q_0\). Cut at \(q_1\) unless
its block is the second member of a pair, in which case cut at
\(q_2\). The cut is an image boundary and has a preceding return
block starting at or after \(N\). This gives \(|H|\le N+10\)
and the exact factorization \(x=H g_i(z_0)\). The decoded word
\(z_0\) is infinite because each decoded letter uses at most two
return blocks.

To verify the legality of \(z_0\), consider any factor \(v\) of
length at most \(4\) in it. Its image \(g_i(v)\) uses at most eight
return blocks. Include the immediately preceding and following
return blocks in \(x\), and append the zero that marks the end
of the following block. The resulting segment has length at most
\(4(8+2)+1=41\). Its start lies at or after \(N\), including when
\(v\) is an initial factor of \(z_0\). By right extension to
length \(41\), the segment occurs in \(s_i\). The zeros identify
all of its return blocks there. The two neighbors guarantee the
same pairing decisions on the core. The occurrence in \(s_i\)
thus decodes to an occurrence of \(v\) in \(t\). All members of
\(F_5\) have length at most \(4\), so \(z_0\) avoids \(F_5\).
A square in \(z_0\) would map to a nonempty square in \(x\),
since \(g_i\) is nonerasing. Hence \(z_0\in\Y\).
\end{proof}

\begin{lemma}[Uniform five-letter desubstitution]\label{lem:five}
Every \(z\in\Y\) has the form \(z=p f(z')\), where \(p\) is a
prefix of length at most \(33\) and \(z'\in\Y\).
\end{lemma}

\begin{proof}
For every position \(s\ge33\) the length-\(8\) factor of \(z\)
starting at \(s\) has \(32\) letters on both sides in \(z\).
Since \(z\in\Y\), \eqref{eq:window} puts the factor in \(\Fac_8(t)\).

In the displayed images of \(f\), \(0,2,4\) occur exactly at image
starts, and \(1,3\) occur only as last letters of two-letter images.
Each \(1\) or \(3\) in \(t=f(t)\) is followed by an image start.
These are local rules on factors of length two, so position \(33\)
or \(34\) in \(z\) is an image start. From that position, read \(01\)
if \(0\) is followed by \(1\), and read the single letter \(0\)
otherwise. Read \(23\) or \(21\) at a \(2\), and read the single
letter \(4\) at a \(4\). These cases follow from the length-two
factors of \(t\) starting at those symbols. Each read ends just
before another image start. Thus the rule reads the whole suffix,
length \(32\) or \(33\) is discarded, and \(z=p f(z')\) for an
infinite word \(z'\).

The morphism \(f\) is nonerasing, so a square in \(z'\) would give
a square in \(z\). Suppose \(v\in F_5\) does occur in \(z'\).
Then \(f(v)\) begins at an image start at or after position \(33\)
in \(z\). For \(v\in F_5\) we have \(|f(v)|\le6\). This follows
from \(|v|\le3\) except at \(\tern{3423}\) and \(\tern{4234}\),
whose images have lengths \(6\) and \(5\), respectively.
Include the first symbol of the next image. The resulting
segment has length at most \(7\), and its right extension to
length \(8\) occurs in \(t\). Both its first symbol and the
appended symbol mark image starts in that occurrence in \(t\).
They fix the decoding, including at a terminal \(0\) of \(f(v)\).
Hence \(v\) occurs in the preimage \(t\) of \(t=f(t)\), a
contradiction. The right boundary matters because, for example,
\(f(\tern{14})=\tern{230}\) is a prefix of
\(f(\tern{10})=\tern{2301}\). In fact \(\tern{10}\) occurs
in \(f^4(0)\), so the bare image \(\tern{230}\) occurs in \(t\)
even though \(\tern{14}\) is forbidden. Therefore \(z'\) avoids
\(F_5\) and belongs to \(\Y\).
\end{proof}

\subsection{Prefix bounds and the arithmetic corollary}\label{sec:arithmetic}

Starting from \(z_0\) in Lemma~\ref{lem:ternary}, apply
Lemma~\ref{lem:five} successively. Its legality conclusion is the
hypothesis needed for the next application. Thus for every \(j\ge0\)
there are \(z_{j+1}\in\Y\) and a prefix \(p_j\) with
\[
z_j=p_j f(z_{j+1}),\qquad |p_j|\le33.
\]
In particular, for every \(k\ge1\) the original ternary word \(x\)
has the factorization
\begin{equation}\label{eq:iteration}
x=H\,g_i(p_0)\,g_i(f(p_1))\cdots
     g_i(f^{k-1}(p_{k-1}))\,g_i(f^k(z_k)).
\end{equation}

Write \(m_k=\min_{a\in\Sigma_5}|f^k(a)|\) and
\(L_k=\max_{a\in\Sigma_5}|f^k(a)|\). In the order \(0,1,2,3,4\),
the third iterates are
\[
(\tern{0123421},\tern{0423},\tern{01},\tern{0421},\tern{0123}).
\]
Every third iterate contains \(0\) and has length at least \(2\).
Since \(f^3(0)\) contains all five letters, every \(f^6(a)\) contains
all five letters. Also \(|f^6(a)|\le 64\). It follows that
\begin{equation}\label{eq:growth}
m_k\ge2^{\lfloor k/3\rfloor},\qquad L_{k+3}\ge2L_k,\qquad
L_k\le64L_{k-6}\le64m_k\quad(k\ge6).
\end{equation}
For the last inequality, use \(f^k(a)=f^{k-6}(f^6(a))\).
The sequence \(L_k\) is nondecreasing because \(f\) is nonerasing.
For \(1\le d\le k\) we have
\(L_{k-d}\le 2^{-\lfloor d/3\rfloor}L_k\), and thus
\begin{equation}\label{eq:sum}
\sum_{j=0}^{k-1} L_j
\le \left(2+3\sum_{h=1}^\infty 2^{-h}\right)L_k=5L_k.
\end{equation}
The part before \(g_i(f^k(z_k))\) in \eqref{eq:iteration} therefore
has length at most \(|H|+8\cdot33\cdot5L_k=|H|+1320L_k\).

Two of the first six letters of \(z_k\) are equal. Their images
under \(g_i\circ f^k\) are identical factors of \(x\). Let their
length be \(n_k\), their starting positions \(u_k<v_k\), and their
distance \(d_k=v_k-u_k\). Each letter image under \(g_i\circ f^k\)
has length between \(2m_k\) and \(8L_k\). There are at most five of
these images between the two starts. By \eqref{eq:growth}, for \(k\ge6\),
\begin{equation}\label{eq:dist}
\begin{aligned}
n_k &\ge2m_k,\qquad
d_k\le40L_k\le2560m_k\le1280n_k,\\
u_k-1 &\le |H|+1360L_k.
\end{aligned}
\end{equation}
The equality of the two factors supplies all comparisons needed
for period \(d_k\) on the length-\((d_k+n_k)\) segment starting at
\(u_k\). If \(d_k\le n_k\), this segment contains a square of root
length \(d_k\) at \(u_k\). Therefore \(d_k>n_k\).

Suppose now that \(x\) is a letter permutation of a suffix of a
ternary digit word \(\mathbf a=a_1a_2\cdots\), after discarding
\(\ell\) digits. The inverse permutation gives identical factors
at positions \(\ell+u_k,\ell+v_k\) in \(\mathbf a\). Let \(U_k\) be
the prefix ending at \(\ell+u_k-1\), and \(V_k\) the segment from
\(\ell+u_k\) to \(\ell+v_k-1\). Define \(V^w\) for \(w>1\) to be the
length-\(\lceil w|V|\rceil\) prefix of \(V^\omega\). Using
\eqref{eq:dist} with \(w=1+1/1280\), we have
\(\lceil d_k/1280\rceil\le n_k\).
The periodic segment in \(\mathbf a\) therefore gives
\begin{equation}\label{eq:abl}
\begin{aligned}
U_kV_k^w &\preceq \mathbf a,\qquad
          |V_k|=d_k>n_k\longrightarrow\infty,\\
\frac{|U_k|}{|V_k|}
&\le \frac{\ell+|H|+1360L_k}{d_k}
 \le \ell+|H|+43520.
\end{aligned}
\end{equation}
Here \(\preceq\) denotes the prefix relation and
\(L_k\le64m_k\le32n_k<32d_k\).
The transfer to \(\mathbf a\) uses only equality and length of
corresponding factors under the permutation.

\begin{proof}[Proof of Theorem T]
If both alternatives fail, choose two distinct members of \(P\)
with finitely many occurrences. Cut the digit word beyond their
last occurrences and beyond every square start. Encode a palindrome
\(aba\) by the arrow \(a\to b\). Under permutations of \(\Sigma_3\),
an unordered pair of distinct arrows has a common target, a common
source, opposite directions on the same edge, or forms a directed
path of length two. Representatives of the four types are
\[
\{\tern{010},\tern{212}\},\quad B_1,\quad B_2,\quad
\{\tern{010},\tern{121}\}.
\]

For \(\{\tern{010},\tern{121}\}\), the jump graph of Shallit and Shur
\cite[\S3.1, Figure 1 and Lemma 21(2)]{SS} excludes an infinite
square-free avoiding word. Every length-six factor of a square-free
word contains a jump \(aba\), since three distinct letters in every
consecutive triple would force a square of length six. Both
forbidden jumps belong to
\(\{\tern{010},\tern{121},\tern{202}\}\), one side of the graph.
With only \(\tern{202}\) left on that side, successive two-edge
returns give label pairs \(11,22\), or \(33\). The lemma excludes
an internal \(11\). A pair \(22\) followed by \(22\) gives the
forbidden \(222\) with a closed two-edge prefix. A pair \(22\)
followed by \(33\) gives an internal \(223\), also forbidden.
A tail of pairs \(33\) gives a forbidden \(333\) with a closed
two-edge prefix. An infinite walk is therefore impossible.

For \(\{\tern{010},\tern{212}\}\), the suffix also avoids
\(\tern{02120}\). The proof of \cite[Theorem 3.1]{AR} supplies
prefixes \(Y_k Z_k c_k Z_k\) on the normalized suffix, with
\(c_k\) a letter, \(|Y_k|\le2^{k+2}\) and \(|Z_k|\ge2^k/3\).
For all sufficiently large \(k\), the period word \(Z_kc_k\)
and the final \(Z_k\) give a \(3/2\)-power. After the inverse
permutation and the restoration of the fixed initial \(\ell\)
digits, the roots tend to infinity and the prefix-to-root ratios
are at most \(\ell+12\). For \(B_1\) and \(B_2\) the required
prefixes are those of \eqref{eq:abl}.

In all three cases with controlled prefixes, take a subsequence
on which the root lengths strictly increase. Since \(\xi\) is
irrational, its original digit word is not eventually periodic.
The criterion of Adamczewski, Bugeaud, and Luca
\cite[Theorem 1]{ABL}, applied in base \(3\) to that word, makes
\(\xi-\lfloor\xi\rfloor\) transcendental. This contradicts
algebraicity of \(\xi\).
\end{proof}

\subsection{Finite verification}\label{sec:computation}

The supplement's \texttt{verify\_window.py} enumerates the words
of length \(72\) on \(\Sigma_5\) that are square-free and avoid
\(F_5\). Starting with the empty word, it appends each letter.
When the prefix is legal, a new forbidden factor or square must
end at the appended letter. The checks are therefore all
forbidden suffixes and every square suffix with root length at
most half the new length. Induction on length establishes
completeness of this enumeration.

The enumeration gives \(1{,}012\) legal words of length \(72\)
and \(33\) distinct centers of length \(8\) at zero-based offsets
\(32\) through \(39\). Each center occurs in \(f^8(0)\), a prefix
of \(t\) of length \(72\). The JSON certificate records the
population at each length, one context witness and one prefix
offset per center, and a SHA-256 digest of the lexicographic stream
of legal length-\(72\) words. Running
\texttt{python verify\_window.py} beside the certificate repeats
the enumeration and checks the recorded data.

This computation verifies \eqref{eq:window}. Legality of the
arbitrary infinite preimages, the iteration, and the prefix
estimates come from the proofs of Lemmas~\ref{lem:ternary}
and~\ref{lem:five} and Section~\ref{sec:arithmetic}.

